\documentclass[a4paper,11pt]{article}

% ──── Required packages (order per project convention) ────
\usepackage{amsmath, amssymb, amsthm}
\usepackage{physics}       % \dd, \grad, \curl, \div, bra-ket
\usepackage{bm}            % \bm{} for bold-italic vectors
\usepackage{siunitx}       % \SI{3e8}{\meter\per\second}
\usepackage{booktabs}      % \toprule \midrule \bottomrule
\usepackage{geometry}
\geometry{margin=25mm}
\usepackage{enumitem}
\usepackage{hyperref}      % LAST package

\title{Theory and Algorithm Notes for \texttt{MieScat\_Py}\\
\large Lorenz--Mie Scattering by Neutral and Charged Homogeneous Spheres}
\author{Technical Note for MieScat\_Py v0.2.0}
\date{\today}

\begin{document}
\maketitle
\tableofcontents
\newpage

%=============================================================================
\section{Introduction}
\label{sec:introduction}
%=============================================================================

\texttt{MieScat\_Py} solves the classical problem of electromagnetic scattering
by a single homogeneous isotropic sphere embedded in a nonabsorbing medium.
The implementation follows the formulations in Bohren and Huffman
\cite{bohren1983} (hereafter BH83) and Mishchenko et al.\ \cite{mishchenko2002},
with an extension to electrically charged spheres based on Klacka and Kocifaj
\cite{klacka2010} (hereafter KK10).  The forward and backward complex
scattering amplitudes used for Complex Amplitude Sensing (CAS) are defined
according to Moteki \cite{moteki2021} and Moteki and Adachi \cite{moteki2024}.

This note documents the mathematical formulations, numerical algorithms, and
their correspondence to the code.  Sec.~\ref{sec:mie-theory} covers the
standard Mie theory for neutral spheres, from the problem definition through to
the computed observables.  Sec.~\ref{sec:charged-sphere} describes the
extension to charged spheres.  Sec.~\ref{sec:cas} defines the CAS-specific
observables.  Sec.~\ref{sec:implementation} summarizes the computational
implementation.


%=============================================================================
\section{Mie theory for neutral spheres}
\label{sec:mie-theory}
%=============================================================================

\subsection{Problem definition and notation}
\label{sec:problem-definition}

Consider a sphere of diameter $d_\mathrm{p}$ (radius
$r_\mathrm{p} = d_\mathrm{p}/2$) with complex refractive index
\begin{equation}
  m_\mathrm{p} = n_\mathrm{p} + i\, \kappa_\mathrm{p}
  \label{eq:mp-def}
\end{equation}
embedded in a nonabsorbing medium with real refractive index $m_\mathrm{m}$.
A monochromatic plane wave of vacuum wavelength $\lambda_0$ illuminates the
sphere.  The vacuum wavenumber, medium wavenumber, size parameter, and relative
refractive index are
\begin{align}
  k_0 &= \frac{2\pi}{\lambda_0},
  &
  k   &= m_\mathrm{m}\, k_0,
  &
  x   &= k\, r_\mathrm{p}
       = \frac{\pi\, d_\mathrm{p}\, m_\mathrm{m}}{\lambda_0},
  &
  m_\mathrm{r} &= \frac{m_\mathrm{p}}{m_\mathrm{m}}.
  \label{eq:basic-params}
\end{align}
The number of terms retained in the partial wave expansion is
\begin{equation}
  N_\mathrm{stop} = \lfloor |x| + 4\,|x|^{1/3} + 2 \rfloor,
  \label{eq:nstop}
\end{equation}
following the criterion of BH83.

\subsection{Ricatti--Bessel functions and logarithmic derivative}
\label{sec:ricatti-bessel}

The Mie solution is expressed in terms of the Ricatti--Bessel functions
\begin{equation}
  \psi_n(z) = z\, j_n(z),
  \qquad
  \chi_n(z) = -z\, y_n(z),
  \qquad
  \xi_n(z) = \psi_n(z) + i\,\chi_n(z) = z\, h_n^{(1)}(z),
  \label{eq:ricatti-bessel-def}
\end{equation}
where $j_n$, $y_n$, and $h_n^{(1)}$ are the spherical Bessel, Neumann, and
Hankel functions of the first kind, respectively.

The logarithmic derivative of $\psi_n$ at argument $y = m_\mathrm{r}\, x$ is
\begin{equation}
  D_n(y) = \frac{\dd}{\dd y}\ln \psi_n(y)
         = \frac{\psi_n'(y)}{\psi_n(y)}.
  \label{eq:log-deriv-def}
\end{equation}
It is computed by the numerically stable \emph{downward} recurrence (BH83,
Sec.~4.8)
\begin{equation}
  D_{n-1}(y) = \frac{n}{y} - \frac{1}{D_n(y) + n/y},
  \qquad n = N_\mathrm{mx},\, N_\mathrm{mx}-1,\, \ldots,\, 1,
  \label{eq:log-deriv-recurrence}
\end{equation}
starting from $D_{N_\mathrm{mx}}(y) = 0$, where
$N_\mathrm{mx} = \lfloor \max(N_\mathrm{stop},\, |y|) + 15 \rfloor$.

$\psi_n(x)$ is computed via the ratio
$R_n = \psi_n(x)/\psi_{n-1}(x)$ using downward recurrence
(Mishchenko et al.\ \cite{mishchenko2002}):
\begin{equation}
  R_n = \frac{1}{(2n+1)/x - R_{n+1}},
  \quad n = N_\mathrm{mx}-1,\, \ldots,\, 0,
  \qquad
  R_{N_\mathrm{mx}} = \frac{x}{2N_\mathrm{mx}+1},
  \label{eq:psi-ratio-recurrence}
\end{equation}
followed by $\psi_0(x) = R_0 \cos x$ and
$\psi_n(x) = R_n\, \psi_{n-1}(x)$ for $n \ge 1$.

$\chi_n(x)$ is computed by upward recurrence, which is stable for this
function:
\begin{equation}
  \chi_0(x) = -\cos x, \qquad
  \chi_1(x) = \frac{\chi_0(x)}{x} - \sin x, \qquad
  \chi_n(x) = \frac{2n-1}{x}\,\chi_{n-1}(x) - \chi_{n-2}(x).
  \label{eq:chi-recurrence}
\end{equation}

\subsection{Partial wave coefficients}
\label{sec:partial-wave}

For a neutral (uncharged) sphere, the partial wave (Mie) coefficients $a_n$ and
$b_n$ are (BH83, Eq.~4.88)
\begin{align}
  a_n &= \frac{
    \bigl(D_n(y)/m_\mathrm{r} + n/x\bigr)\,\psi_n(x) - \psi_{n-1}(x)
  }{
    \bigl(D_n(y)/m_\mathrm{r} + n/x\bigr)\,\xi_n(x) - \xi_{n-1}(x)
  },
  \label{eq:an-neutral} \\[6pt]
  b_n &= \frac{
    \bigl(m_\mathrm{r}\,D_n(y) + n/x\bigr)\,\psi_n(x) - \psi_{n-1}(x)
  }{
    \bigl(m_\mathrm{r}\,D_n(y) + n/x\bigr)\,\xi_n(x) - \xi_{n-1}(x)
  },
  \label{eq:bn-neutral}
\end{align}
where $y = m_\mathrm{r}\, x$ and $n = 1, 2, \ldots, N_\mathrm{stop}$.
The coefficients $a_n$ correspond to the electric (TM) multipoles and $b_n$ to
the magnetic (TE) multipoles.

\subsection{Optical efficiency factors}
\label{sec:efficiency-factors}

Given $a_n$ and $b_n$, the optical efficiency factors are (BH83,
Eqs.~4.61--4.62)
\begin{align}
  Q_\mathrm{sca} &= \frac{2}{x^2}
    \sum_{n=1}^{N_\mathrm{stop}} (2n+1)\bigl(|a_n|^2 + |b_n|^2\bigr),
  \label{eq:qsca} \\
  Q_\mathrm{ext} &= \frac{2}{x^2}
    \sum_{n=1}^{N_\mathrm{stop}} (2n+1)\,\Re(a_n + b_n),
  \label{eq:qext} \\
  Q_\mathrm{abs} &= Q_\mathrm{ext} - Q_\mathrm{sca},
  \label{eq:qabs} \\
  Q_\mathrm{back} &= \frac{1}{x^2}
    \left|\sum_{n=1}^{N_\mathrm{stop}} (2n+1)(-1)^n (a_n - b_n)\right|^2.
  \label{eq:qback}
\end{align}
These are dimensionless ratios of the respective cross sections to the
geometric cross section $\pi r_\mathrm{p}^2$.

\subsection{Scattering amplitudes}
\label{sec:scattering-amplitudes}

The angle-dependent functions $\pi_n(\cos\theta)$ and $\tau_n(\cos\theta)$
are defined as
\begin{equation}
  \pi_n(\cos\theta) = \frac{P_n^1(\cos\theta)}{\sin\theta}, \qquad
  \tau_n(\cos\theta) = \frac{\dd P_n^1(\cos\theta)}{\dd\theta},
  \label{eq:pi-tau-def}
\end{equation}
where $P_n^1$ is the associated Legendre function of degree $n$ and order~1.
They satisfy the recurrence relations (with $\mu = \cos\theta$,
$\pi_1 = 1$, $\pi_2 = 3\mu$)
\begin{equation}
  \pi_n = \frac{2n-1}{n-1}\,\mu\,\pi_{n-1} - \frac{n}{n-1}\,\pi_{n-2},
  \qquad
  \tau_n = n\,\mu\,\pi_n - (n+1)\,\pi_{n-1}.
  \label{eq:pi-tau-recurrence}
\end{equation}

The elements of the $2 \times 2$ amplitude scattering matrix (BH83, Eq.~4.74)
are
\begin{align}
  S_1(\theta) &= \sum_{n=1}^{N_\mathrm{stop}}
    \frac{2n+1}{n(n+1)}\bigl(a_n\,\pi_n + b_n\,\tau_n\bigr),
  \label{eq:s1} \\
  S_2(\theta) &= \sum_{n=1}^{N_\mathrm{stop}}
    \frac{2n+1}{n(n+1)}\bigl(a_n\,\tau_n + b_n\,\pi_n\bigr).
  \label{eq:s2}
\end{align}
The code converts these to the amplitude matrix convention of
Mishchenko et al.\ \cite{mishchenko2002}:
\begin{equation}
  S_{11}(\theta) = \frac{S_2(\theta)}{-ik}, \qquad
  S_{22}(\theta) = \frac{S_1(\theta)}{-ik}.
  \label{eq:s11-s22}
\end{equation}
Note the index swap: $S_{11}$ is derived from $S_2$, and $S_{22}$ from $S_1$.
Both have the dimension of length.

\subsection{Mass-normalized cross sections}
\label{sec:mass-cross-sections}

The geometric cross section, particle volume, and particle mass are
\begin{equation}
  C_\mathrm{geo} = \frac{\pi d_\mathrm{p}^2}{4},
  \qquad
  V_\mathrm{p} = \frac{\pi d_\mathrm{p}^3}{6},
  \qquad
  M_\mathrm{p} = \rho_\mathrm{p}\, V_\mathrm{p},
  \label{eq:mass-def}
\end{equation}
where $\rho_\mathrm{p}$ is the particle material density.
The mass-normalized cross sections are
\begin{equation}
  \mathrm{MSC} = \frac{C_\mathrm{geo}\, Q_\mathrm{sca}}{M_\mathrm{p}},
  \qquad
  \mathrm{MEC} = \frac{C_\mathrm{geo}\, Q_\mathrm{ext}}{M_\mathrm{p}},
  \qquad
  \mathrm{MAC} = \frac{C_\mathrm{geo}\, Q_\mathrm{abs}}{M_\mathrm{p}}.
  \label{eq:mass-cross-sections}
\end{equation}
When $d_\mathrm{p}$ is given in \si{\meter} and $\rho_\mathrm{p}$ in
\si{\gram\per\cubic\centi\meter}, the code internally converts
$V_\mathrm{p}$ to \si{\cubic\centi\meter} so that
$M_\mathrm{p}$ is in \si{\gram} and the mass-normalized cross sections are in
\si{\meter\squared\per\gram}.


%=============================================================================
\section{Extension to charged spheres}
\label{sec:charged-sphere}
%=============================================================================

For a sphere carrying $n_\mathrm{e}$ elementary charges $e$ in vacuum
($m_\mathrm{m} = 1$), the formulation of KK10 introduces a
charge-dependent $G$-parameter that modifies $a_n$ and $b_n$.  The total charge
is $Q = n_\mathrm{e}\, e$ and the electrostatic surface potential is
\begin{equation}
  V_\mathrm{surf} = \frac{Q}{4\pi\epsilon_0\, r_\mathrm{p}}.
  \label{eq:surf-potential}
\end{equation}
The quantum scattering rate is estimated as (KK10, Eq.~37)
\begin{equation}
  \gamma_\mathrm{s} = f_\gamma \,\frac{k_\mathrm{B} T}{\hbar},
  \label{eq:gamma-s}
\end{equation}
where $f_\gamma$ is a dimensionless correction factor, $k_\mathrm{B}$ is the
Boltzmann constant, $T = \SI{298}{\kelvin}$, and $\hbar$ is the reduced Planck
constant.  The dimensionless charge parameter is (KK10, Eq.~26)
\begin{equation}
  G = \frac{e\, V_\mathrm{surf}}{m_\mathrm{e}\, c^2}
      \cdot \frac{1}{x}
      \cdot \frac{1 + i\,(\gamma_\mathrm{s}/\omega)}
                 {1 + (\gamma_\mathrm{s}/\omega)^2},
  \label{eq:g-parameter}
\end{equation}
where $m_\mathrm{e}$ is the electron mass, $c$ is the speed of light, and
$\omega = k\, c$ is the angular frequency.

The modified partial wave coefficients for the charged sphere are
\begin{align}
  a_n &= \frac{
    \bigl[(1 + nG/x)\, D_n/m_\mathrm{p} + n/x\bigr]\,\psi_n
    - (1 + G\, D_n/m_\mathrm{p})\,\psi_{n-1}
  }{
    \bigl[(1 + nG/x)\, D_n/m_\mathrm{p} + n/x\bigr]\,\xi_n
    - (1 + G\, D_n/m_\mathrm{p})\,\xi_{n-1}
  },
  \label{eq:an-charged} \\[6pt]
  b_n &= \frac{
    \bigl[m_\mathrm{p}\, D_n + (n/x)(1 - Gx/n)\bigr]\,\psi_n - \psi_{n-1}
  }{
    \bigl[m_\mathrm{p}\, D_n + (n/x)(1 - Gx/n)\bigr]\,\xi_n - \xi_{n-1}
  },
  \label{eq:bn-charged}
\end{align}
where all Ricatti--Bessel functions are evaluated at~$x$, and
$D_n = D_n(m_\mathrm{p}\, x)$ with $m_\mathrm{r} = m_\mathrm{p}$ (since the
medium is vacuum).  When $G = 0$ (i.e., $n_\mathrm{e} = 0$),
\eqref{eq:an-charged}--\eqref{eq:bn-charged} reduce to the standard Mie
expressions \eqref{eq:an-neutral}--\eqref{eq:bn-neutral} with
$m_\mathrm{r} = m_\mathrm{p}$.  The efficiency factors, scattering amplitudes,
and mass-normalized cross sections are then computed identically to
Sec.~\ref{sec:efficiency-factors}--\ref{sec:mass-cross-sections}.


%=============================================================================
\section{Complex Amplitude Sensing observables}
\label{sec:cas}
%=============================================================================

For the CAS protocol \cite{moteki2021,moteki2024}, the relevant observables
are the forward scattering amplitude $S(0)$ and a backward scattering
amplitude $S_\mathrm{bak}$.

At $\theta = 0$ the angular functions satisfy $\pi_n(1) = \tau_n(1) = 1$ for
all~$n$, so $S_1(0) = S_2(0)$ and hence $S_{11}(0) = S_{22}(0)$.  The forward
scattering amplitude is
\begin{equation}
  S(0) = \frac{S_{11}(0) + S_{22}(0)}{2} = S_{11}(0).
  \label{eq:s-fwd}
\end{equation}
At $\theta = \pi$ the angular functions satisfy $\pi_n(-1) = (-1)^{n+1}$ and
$\tau_n(-1) = (-1)^n$, giving $S_1(\pi) = -S_2(\pi)$.  The backward
scattering amplitude for CAS-v2 is
\begin{equation}
  S_\mathrm{bak} = \frac{-S_{11}(\pi) + S_{22}(\pi)}{\sqrt{2}}.
  \label{eq:s-bak}
\end{equation}
The code returns $\Re[S(0)]$, $\Im[S(0)]$, $\Re[S_\mathrm{bak}]$, and
$\Im[S_\mathrm{bak}]$, all having the dimension of length.


%=============================================================================
\section{Computational implementation}
\label{sec:implementation}
%=============================================================================

\begin{table}[htbp]
  \caption{Modules in \texttt{MieScat\_Py} v0.2.0 and their roles.}
  \label{tab:modules}
  \centering
  \begin{tabular}{@{}ll@{}}
    \toprule
    Module & Role \\
    \midrule
    \texttt{\_mie\_core.py}            & Shared helper functions
      (\eqref{eq:log-deriv-recurrence}--\eqref{eq:s11-s22}) \\
    \texttt{miescat.py}                & Neutral sphere:
      $Q$-factors, $S_{11}$, $S_{22}$, MSC/MEC/MAC \\
    \texttt{miescat\_charged.py}       & Charged sphere (KK10):
      adds $V_\mathrm{surf}$, $G$-parameter \\
    \texttt{mie\_complex\_amplitudes.py} & CAS observables:
      $S(0)$ and $S_\mathrm{bak}$ \\
    \bottomrule
  \end{tabular}
\end{table}

The key numerical considerations are as follows.
The logarithmic derivative $D_n(y)$ is computed by \emph{downward} recurrence
\eqref{eq:log-deriv-recurrence}, which is numerically stable for complex~$y$
(BH83, Sec.~4.8).
$\psi_n(x)$ is computed via the ratio method
\eqref{eq:psi-ratio-recurrence} using downward recurrence on the ratios $R_n$,
which avoids the catastrophic cancellation that plagues upward recurrence for
large~$x$ (Mishchenko et al.\ \cite{mishchenko2002}).
$\chi_n(x)$ is computed by upward recurrence \eqref{eq:chi-recurrence}, which
is stable for this function.
The recurrence relations for $D_n$, $R_n$, $\psi_n$, $\chi_n$, $\pi_n$,
and $\tau_n$ remain as sequential loops because each step depends on the
previous value.

In v0.2.0, the following element-wise Python loops were replaced by NumPy
vectorized operations:
(i)~the weight arrays $(2n+1)$, $(2n+1)/[n(n+1)]$, and $(-1)^n$;
(ii)~the partial wave coefficients $a_n$, $b_n$
(\eqref{eq:an-neutral}--\eqref{eq:bn-neutral} or
 \eqref{eq:an-charged}--\eqref{eq:bn-charged}) for all $n$ simultaneously;
(iii)~the scattering amplitude sums \eqref{eq:s1}--\eqref{eq:s2} over $n$
for all angles $\theta$ simultaneously via array broadcasting.

\begin{table}[htbp]
  \caption{Correspondence between mathematical symbols and code variables.}
  \label{tab:variables}
  \centering
  \begin{tabular}{@{}lll@{}}
    \toprule
    Symbol & Code variable & Description \\
    \midrule
    $\lambda_0$        & \texttt{wl\_0}      & Vacuum wavelength \\
    $m_\mathrm{m}$     & \texttt{m\_m}       & Medium refractive index \\
    $d_\mathrm{p}$     & \texttt{d\_p}       & Particle diameter \\
    $m_\mathrm{p}$     & \texttt{m\_p}       & Complex particle refractive index \\
    $m_\mathrm{r}$     & \texttt{m\_r}       & Relative refractive index \\
    $x$                & \texttt{x}          & Size parameter \\
    $N_\mathrm{stop}$  & \texttt{nstop}      & Number of expansion terms \\
    $D_n(y)$           & \texttt{DD[n]}      & Logarithmic derivative \\
    $\psi_n(x)$        & \texttt{psi[n]}     & Ricatti--Bessel (1st kind) \\
    $\chi_n(x)$        & \texttt{chi[n]}     & Ricatti--Bessel (2nd kind) \\
    $\xi_n(x)$         & \texttt{xi[n]}      & Ricatti--Bessel (3rd kind) \\
    $a_n$              & \texttt{a[n]}       & Electric multipole coefficient \\
    $b_n$              & \texttt{b[n]}       & Magnetic multipole coefficient \\
    $\pi_n$            & \texttt{pi\_n[n,:]} & Angular function \\
    $\tau_n$           & \texttt{tau\_n[n,:]}& Angular function \\
    $S_1$, $S_2$       & \texttt{S1}, \texttt{S2}  & BH83 amplitude elements \\
    $S_{11}$, $S_{22}$ & \texttt{S11}, \texttt{S22} & Mishchenko amplitude elements \\
    $G$                & \texttt{g}          & Charge parameter (KK10) \\
    $V_\mathrm{surf}$  & \texttt{surf\_potential} & Surface potential \\
    $\rho_\mathrm{p}$  & \texttt{dens\_p}    & Particle density \\
    \bottomrule
  \end{tabular}
\end{table}


\clearpage
\begin{thebibliography}{9}
\bibitem{bohren1983}
  C.~F. Bohren and D.~R. Huffman,
  \textit{Absorption and Scattering of Light by Small Particles},
  John Wiley \& Sons, 1983.

\bibitem{mishchenko2002}
  M.~I. Mishchenko, L.~D. Travis, and A.~A. Lacis,
  \textit{Scattering, Absorption, and Emission of Light by Small Particles},
  Cambridge University Press, 2002.

\bibitem{klacka2010}
  J.~Klacka and M.~Kocifaj,
  ``On the scattering of electromagnetic waves by a charged sphere,''
  \textit{Progress In Electromagnetics Research}, vol.~109, pp.~17--35, 2010.

\bibitem{moteki2021}
  N.~Moteki,
  ``Measuring the complex forward-scattering amplitude of single particles
  by self-reference interferometry: CAS-v1 protocol,''
  \textit{Optics Express}, vol.~29, no.~13, pp.~20688--20714, 2021.
  \href{https://doi.org/10.1364/OE.423175}{doi:10.1364/OE.423175}.

\bibitem{moteki2024}
  N.~Moteki and K.~Adachi,
  ``Measuring the polarized complex forward-scattering amplitudes of single
  particles in unbounded fluid flow: CAS-v2 protocol,''
  \textit{Optics Express}, vol.~32, no.~21, pp.~36500--36522, 2024.
  \href{https://doi.org/10.1364/OE.533776}{doi:10.1364/OE.533776}.
\end{thebibliography}

\bigskip
\noindent\textbf{Acknowledgment.}\quad
This document was prepared with the assistance of Claude (Anthropic).
The author assumes full responsibility for the content.

\end{document}
